\begin{abstract}
The Transformer architecture (Vaswani et al., 2017) has become the backbone of modern natural language processing, yet most open-source implementations rely on high-level PyTorch modules that abstract away the underlying mathematical machinery. We present a complete, from-scratch PyTorch implementation of the Transformer that reproduces every component---MultiHeadAttention, PositionwiseFeedForward, and Positional Encoding---using only primitive tensor operations, without \texttt{nn.MultiheadAttention}, \texttt{nn.Transformer}, or \texttt{torchtext}.

Our implementation is validated empirically on the \textbf{Tatoeba English $\rightarrow$ German} sentence-pair collection (CC-BY release via \url{https://www.manythings.org/anki/}; $\sim$330k pairs, of which $\sim$324k are used for training). With a compact configuration ($d_{\text{model}} = 256$, 4 encoder and 4 decoder layers, 8 heads, $d_{\text{ff}} = 1024$, label smoothing $\epsilon = 0.1$, batch size 64, warmup 4{,}000 steps), SentencePiece BPE tokenization with an 8k-piece vocabulary per language, and pre-layer normalization ($\approx$13.5 M parameters), the model trains for 25 epochs in $\approx$ 2.5 hours on a single Google Colab GPU (NVIDIA T4) and reaches a corpus-level \textbf{BLEU = 44.00} greedy (\textbf{45.30} with beam search, beam $=4$) and \textbf{chrF2 = 62.41} (\textbf{63.37} with beam search) on the test set ($n = 3{,}312$) at validation perplexity \textbf{9.29}. A suite of 17 unit tests serves as a formal correctness contract for every architectural component.

This result was reached through an explicit diagnosis-and-correction process that we report in full: an initial word-level control run (Post-LN, batch 32, warmup 2000), stopped after 8 epochs, collapsed to BLEU $= 0.38$ at validation perplexity 149.84. A structured analysis showed this failure to be one of sub-training rather than architecture---the validation loss was still descending and all masking and shape contracts held---and it motivated three corrections adopted by the final configuration: pre-layer normalization, a corrected learning-rate warmup schedule, and subword (BPE) tokenization.

We further provide a self-contained data pipeline supporting both word-level and subword (SentencePiece BPE) tokenization, a built-from-scratch vocabulary, and length-bucketed batching. Because attention weights are exposed as first-class outputs of the forward pass, the implementation also supports interpretability analysis: we report per-sentence heatmaps of encoder self-attention, decoder self-attention, and decoder cross-attention on the final layer (Section~5.5), where the under-trained control produces diffuse cross-attention maps whereas the final checkpoint exhibits sharp source--target alignment peaks on short test sentences.

Future work includes multi-seed replication, a baseline comparison with fully matched training schedules, the remaining controlled ablations (label smoothing, warmup, model depth/width), and an extension of the attention analysis to longer sequences. All code, hyperparameters, and experimental logs are publicly available, designed for transparency, auditability, and pedagogical use.
\end{abstract}
